

Theorem~\ref{thm:mop-functional-forms} follows immediately as a consequence of the following results, Lemmas \ref{lem:mop-1-formula} and \ref{lem:mop-0-formula}.
Here, an $A$ superscript is used to denote ancestral trajectories, so that $x_{1:n,j}^{A, P,\theta}$ contains the history of the prediction particle $x_{n,j}^{P,\theta}$ at observation times $\seq{1}{n}$.
Similarly, $g_{i,j}^{A,P,\theta}$ is the ancestral weight of $x_{n,j}^{P,\theta}$ at observation time $i$.

\begin{lem}
\label{lem:mop-1-formula}
Write $\nabla_\theta \hat\ell^\alpha(\theta)$ for the DMOP-$\alpha$ gradient estimate (or, equivalently, the gradient of MOP-$\alpha$ when $\theta=\phi$).
Consider the case where we use the after-resampling conditional likelihood estimate so that $\hat\lik(\theta) = \prod_{n=1}^N L_n^{A, \theta, \alpha}$.
When $\alpha=1$,
    \begin{equation}
        \nabla_\theta \hat{\ell}^1(\theta) 
        = \frac{1}{J}\sum_{j=1}^J \nabla_\theta \log f_{Y_{1:N}|X_{1:N}}\big(y_{1:N}^* | x_{1:n,j}^{A, F,\theta};\theta\big),
    \end{equation}
    yielding the bootstrap filter estimator of \cite{poyiadjis11} and \cite{scibior21}.
\end{lem}

\begin{proof}[Proof.]
    Consider the case of MOP-$\alpha$ when $\alpha=1$ and $\theta=\phi$. We then have a nice telescoping product property for the after-resampling likelihood estimate:
\begin{eqnarray}
    \hat{\lik}^1(\theta) &:=& \prod_{n=1}^N L_n^{A, \theta, \alpha} = \prod_{n=1}^N L_n^\phi \cdot \frac{\sum_{j=1}^J w_{n,j}^{F,\theta}}{\sum_{j=1}^J w_{n,j}^{P,\theta}}= \prod_{n=1}^N L_n^\phi \cdot \frac{\sum_{j=1}^J w_{n,j}^{F,\theta}}{\sum_{j=1}^J w_{n-1,j}^{F,\theta}} \nonumber
\\
&=& \left(\frac{1}{J}\sum_{j=1}^J w_{N,j}^{F,\theta}\right) \prod_{n=1}^N L_n^\phi,
\end{eqnarray}
where the third equality follows from the choice of $\alpha=1$, and the fourth equality is the resulting telescoping property. 
The log-derivative identity lets us decompose the score estimate as
\begin{equation}\label{eq:log-derivative-identity}
\nabla_\theta \hat\ell^1(\theta) = \frac{\nabla_\theta \hat\lik^1(\theta)}{\hat\lik^1(\theta)} = \frac{\nabla_\theta\left(\frac{1}{J}\sum_{j=1}^J w_{N,j}^{F,\theta}\right) \prod_{n=1}^N L_n^\phi}{\prod_{n=1}^N L_n^\phi} =  \frac{1}{J}\sum_{j=1}^J \nabla_\theta w_{N,j}^{F,\theta}.
\end{equation}
From (\ref{eq:log-derivative-identity}), we see that the derivative of the log-likelihood estimate is
\begin{equation}\label{eq:eq:log-derivative2}
    \nabla_\theta \hat{\ell}^1(\theta) := \frac{1}{J}\sum_{j=1}^J \nabla_\theta w_{N,j}^{F,\theta}.
\end{equation}
We proceed to decompose (\ref{eq:eq:log-derivative2}).
First, observe that as $\alpha=1$,
\begin{equation}\label{eq:lemma-s1-ratio}
w_{n,j}^{P,\theta} = w_{n-1,j}^{F,\theta}\frac{g_{n,j}^\theta}{g_{n,j}^\phi} = \prod_{i=1}^n \frac{g_{i,j}^{A,P,\theta}}{g_{i,j}^{A,P,\phi}},
\end{equation}
Note that the right hand side of (\ref{eq:lemma-s1-ratio}) is the cumulative product of measurement density ratios over the ancestral trajectory of the $j$-th prediction particle at timestep $n$.
We then use the log-derivative identity again, yielding the following expression for the gradient of the log-weights as the sum of the log measurement densities over the ancestral trajectory,
\begin{eqnarray}
 \frac{\nabla_\theta w_{n,j}^{P,\theta}}{w_{n,j}^{P,\theta}} = \nabla_\theta \log w_{n,j}^{P,\theta} &=& \nabla_\theta \log \left(\prod_{i=1}^n \frac{g_{i,j}^{A,P,\theta}}{g_{i,j}^{A,P,\phi}}\right) 
 \\
 &=& \nabla_\theta \sum_{i=1}^n \left(\log g_{i,j}^{A,P,\theta} - \log g_{i,j}^{A,P,\phi}\right)
 \\
 &=& \sum_{i=1}^n \nabla_\theta \log g_{i,j}^{A,P,\theta}.
\end{eqnarray}
This is equal to the gradient of the logarithm of the conditional density of the observed measurements given the ancestral trajectory of the $j$-th prediction particle up to timestep $n$,
\begin{eqnarray}  
\nabla_\theta \sum_{n=1}^N \log g_{n,j}^{A,P,\theta} &=& \nabla_\theta \log\left(\prod_{n=1}^N g_{n,j}^{A,P,\theta}\right) 
\\
&=&  \nabla_\theta \log\left(\prod_{n=1}^N f_{Y_n|X_n}\big(y_n^* | x_{n,j}^{A, P,\theta};\theta\big)\right)
\\
&=& \nabla_\theta \log f_{Y_{1:N}|X_{1:N}}\big(y_{1:N}^* | x_{1:n,j}^{A, P,\theta};\theta\big).
\end{eqnarray}
Multiplying both sides of the expression by $w_{N,j}^{P,\theta} $ yields an expression for the gradient of the weights at timestep $N$,
\begin{equation}
\nabla_\theta w_{N,j}^{P,\theta} = w_{N,j}^{P,\theta} \sum_{n=1}^N \nabla_\theta \log g_{n,j}^{A,P,\theta} = w_{N,j}^{P,\theta} \nabla_\theta \log f_{Y_{1:N}|X_{1:N}}\big(y_{1:N}^* | x_{1:n,j}^{A, P,\theta};\theta\big).    
\end{equation}
Substituting the above identity into the log-likelihood decomposition obtained earlier in Equation \ref{eq:log-derivative-identity} yields
\begin{eqnarray}
    \nabla_\theta \hat{\ell}^1(\theta) &=& \frac{1}{J}\sum_{j=1}^J \nabla_\theta w_{N,j}^{F,\theta} =\frac{1}{J}\sum_{j=1}^J \nabla_\theta w_{N,k_j}^{P,\theta}
    \nonumber \\
    &=& \frac{1}{J}\sum_{j=1}^J w_{N,k_j}^{P,\theta} \nabla_\theta \log f_{Y_{1:N}|X_{1:N}}\big(y_{1:N}^* | x_{1:n,k_j}^{A, P,\theta};\theta \big).
\end{eqnarray}
Finally, observing that $\theta=\phi$ implies $w_{N,j}^{F,\theta}=1$, we obtain 
\begin{equation}
    \nabla_\theta \hat{\ell}^1(\theta) = \frac{1}{J}\sum_{j=1}^J \nabla_\theta \log f_{Y_{1:N}|X_{1:N}}\big(y_{1:N}^* | x_{1:n,j}^{A, F,\theta};\theta\big).
\end{equation}
This yields the gradient estimators of \cite{poyiadjis11, scibior21} when applied to the bootstrap filter. 
\end{proof}

Note that the variance of the DMOP-$\alpha$ log-likelihood derivative estimate scales poorly with $N$ when $\theta\neq\phi$. 
This can be seen by observing that
\begin{eqnarray}
    \nabla_\theta \hat{\ell}^1(\theta) &=& \frac{1}{J}\sum_{j=1}^J \nabla_\theta w_{N,j}^{F,\theta} =\frac{1}{J}\sum_{j=1}^J \nabla_\theta w_{N,k_j}^{P,\theta}
\nonumber \\
&=& \frac{1}{J}\sum_{j=1}^J w_{N,k_j}^{P,\theta} \nabla_\theta \log f_{Y_{1:N}|X_{1:N}}\big(y_{1:N}^* | x_{1:n,k_j}^{A, P,\theta}; \theta \big).
\end{eqnarray}
When $\theta\neq\phi$, we see that $w_{N,k_j}^{P,\theta} = O(c^N)$. 
When $\theta=\phi$, this is a special case of the \cite{poyiadjis11} estimator, which has $O(N^4)$ variance by a property of functionals of the particle filter (Lemma 2 of \cite{karjalainen23} or Proposition 11.3 of \citet{chopin20}). 

\begin{lem}
\label{lem:mop-0-formula}
Write $\nabla_\theta \hat\ell^\alpha(\theta)$ for the DMOP-$\alpha$ gradient estimate (or, equivalently, the gradient of MOP-$\alpha$ when $\theta=\phi$).
Consider the case where we use the after-resampling conditional likelihood estimate so that $\hat\lik(\theta) = \prod_{n=1}^N L_n^{A, \theta, \alpha}$.
When $\alpha=0$,
    \begin{equation}
        \nabla_\theta \hat\ell^0(\theta) 
        = \frac{1}{J} \sum_{n=1}^N \sum_{j=1}^J \nabla_\theta \log\left(f_{Y_n|X_{n}}\big(y_n^*|x_{n,j}^{F, \theta}; \theta\big)\right),
    \end{equation}
    yielding the estimate of \cite{naesseth18} when applied to the bootstrap filter. 
\end{lem}
\begin{proof}[Proof.]
First, write
\begin{equation}
s_{n,j} = \frac{f_{Y_n|X_n}\big(y_n^*|x_{n,j}^{P, \theta};\theta\big)}{f_{Y_n|X_n}\big(y_n^*|x_{n,j}^{P, \phi};\theta\big)}
\end{equation}
as shorthand for the measurement density ratios. 
Observe that, when $\alpha=0,$, the likelihood estimate becomes
\begin{eqnarray}
% \nonumber
    \hat{\lik}^0(\theta) := \prod_{n=1}^N L_n^{A, \theta, \alpha} &=& \prod_{n=1}^N L_n^\phi \cdot \frac{\sum_{j=1}^J w_{n,j}^{F,\theta}}{\sum_{j=1}^J w_{n,j}^{P,\theta}} 
    \\
    &=& \prod_{n=1}^N L_n^\phi \cdot \frac{1}{J}\sum_{j=1}^J s_{n,j} 
    \\
    &=& \prod_{n=1}^N L_n^\phi \cdot \frac{1}{J}\sum_{j=1}^J \frac{f_{Y_n|X_n}\big(y_n^*|x_{n,j}^{P, \theta};\theta\big)}{f_{Y_n|X_n}\big(y_n^*|x_{n,j}^{P, \phi};\theta\big)}.
\end{eqnarray}
We lose the nice telescoping property observed in the MOP-$1$ case, but this expression still yields something useful. 
This is because its gradient when $\theta=\phi$ is therefore 
\begin{align}
    \nabla_\theta \hat{\ell}^0(\theta) &:= \sum_{n=1}^N \nabla_\theta \log\left(L_n^\phi \frac{1}{J} \sum_{j=1}^J s_{n,j}\right) \\
    &= \sum_{n=1}^N \frac{\nabla_\theta \left(L_n^\phi \frac{1}{J} \sum_{j=1}^J s_{n,j}\right)}{\left(L_n^\phi \frac{1}{J} \sum_{j=1}^J s_{n,j}\right)} \\
    &= \sum_{n=1}^N \frac{\sum_{j=1}^J \nabla_\theta s_{n,j}}{\sum_{j=1}^J s_{n,j}} \\
    &= \sum_{n=1}^N \frac{1}{J} \sum_{j=1}^J \frac{\nabla_\theta f_{Y_n|X_{n}}\big(y_n^*|x_{n,j}^{F, \theta}; \theta\big)}{f_{Y_n|X_{n}}(y_n^*|x_{n,j}^{F, \phi}; \phi)} \\
    &= \frac{1}{J} \sum_{n=1}^N \sum_{j=1}^J \nabla_\theta \log\left(f_{Y_n|X_{n}}\big(y_n^*|x_{n,j}^{F, \theta}; \theta\big)\right),
\end{align}
where we use the log-derivative trick in the second equality, observe that $\sum_{j=1}^J s_{n,j} = J$ when $\theta=\phi$ in the fourth equality, and use the log-derivative trick again while noting that $\theta=\phi$ in the fifth equality. This yields the desired result.
\end{proof}



